\documentclass[12pt,a4paper]{article}
\usepackage[a4paper,top=15mm,bottom=15mm,left=18mm,right=18mm]{geometry}
\usepackage[T1]{fontenc}
\usepackage[utf8]{inputenc}
\usepackage{amsmath}
\usepackage{amssymb}
\usepackage{enumitem}
\usepackage{textcomp}

\setlength{\parindent}{0pt}
\setlist[enumerate,1]{label=\textbf{Q\arabic*.}, leftmargin=*, itemsep=6pt, topsep=4pt}
\setlist[enumerate,2]{label=(\alph*), leftmargin=*, itemsep=3pt, topsep=2pt}

\begin{document}

\begin{center}
{\LARGE\bfseries Answer Key: Currency Conversion Graphs}
\end{center}
\vspace{6pt}

\section*{Part 1: Which line shows which currency?}

\begin{enumerate}
\item
\begin{enumerate}
\item Line A shows Australian dollars.
\item A: Australian dollars; B: zloty; C: euro.
\item \(14\) Australian dollars.
\item \pounds5.
\end{enumerate}

\item
\begin{enumerate}
\item Steepest to least steep: Indian rupees, Mexican pesos, Norwegian kroner, Hong Kong dollars. All pass through \((0,0)\).
\item At \pounds0, \(0\) of any currency; every line passes through \((0,0)\).
\end{enumerate}

\item
\begin{enumerate}
\item P: US dollars; Q: Canadian dollars; R: Swiss francs.
\item Singapore dollars: \(y=1.7x\), between P and Q. New Zealand dollars: \(y=2.2x\), steeper than Q.
\item Sam is wrong. Q is not New Zealand dollars; New Zealand dollars rate \(2.20>1.80\), so its line would be steeper than Q.
\end{enumerate}
\end{enumerate}

\section*{Part 2: Measuring steepness with the gradient}

\begin{enumerate}[resume]
\item Gradient \(=24\). Equation \(y=24x\). \pounds1 buys \(24\) South African rand.

\item
\begin{enumerate}
\item S: Japanese yen; T: Indian rupees.
\item S: \(y=200x\); T: \(y=375x\).
\item \(50\,000\) yen.
\item \pounds80.
\end{enumerate}

\item
\begin{enumerate}
\item zloty: \(0,10,25,50\); pounds: \(0,2,5,10\).
\item Straight line from \((0,0)\) to \((50,10)\).
\item Gradient \(=0.2\).
\item \(1\) zloty is worth \pounds0.20, i.e. \(20\) pence.
\end{enumerate}
\end{enumerate}

\section*{Part 3: Flat fees and the y-intercept}

\begin{enumerate}[resume]
\item \(y=1.3(x-10)=1.3x-13\). Gradient \(1.3\); y-intercept \(-13\).
\begin{enumerate}
\item \(\$117\).
\item \pounds110.
\end{enumerate}

\item
\begin{enumerate}
\item Airport M, High street L, Online N.
\item L and N are parallel (same gradient \(1.2\)); different y-intercepts from different flat fees.
\item L: \(y=1.2x-6\); M: \(y=1.1x\); N: \(y=1.2x-18\).
\item Cross at \pounds60, \texteuro66. For amounts less than \pounds60, airport gives more; for amounts more than \pounds60, high street gives more.
\end{enumerate}

\item
\begin{enumerate}
\item Rate: C\$1.80 per \pounds1.
\item Flat fee: \pounds5.
\item \(x=5\) gives \(y=0\); no Canadian dollars.
\item \(y=1.8x-18\) (or \(1.8(x-10)\)). Parallel; \(9\) Canadian dollars lower at every \(x\).
\end{enumerate}
\end{enumerate}

\section*{Part 4: When do two bureaux give the same amount?}

\begin{enumerate}[resume]
\item
\begin{enumerate}
\item \pounds30: airport \texteuro33, high street \texteuro30; airport gives more.
\item \pounds200: airport \texteuro220, high street \texteuro234; high street gives \texteuro14 more.
\item High street gives \(0.10\) euros more per extra \pounds1. After \pounds60, high street is ahead by \(0.1(x-60)\) euros; at \pounds200 this is \texteuro14.
\end{enumerate}

\item Airport: \(y=4.5x\). Town: \(y=5x-15\). Cross at \((30,135)\).
\begin{enumerate}
\item Draw both lines; mark \((30,135)\).
\item \pounds20: airport \(90\) zloty, town \(85\) zloty; airport gives more.
\item \pounds50: airport \(225\) zloty, town \(235\) zloty; town gives more.
\end{enumerate}

\item
\begin{enumerate}
\item Airport: \(y=1.15x\). Online: \(y=1.25x-10\).
\item Same at \pounds100, \(\$115\).
\item \pounds400: airport \(\$460\), online \(\$490\); online gives \(\$30\) more. Mia should use online.
\end{enumerate}

\item
\begin{enumerate}
\item E: \(y=190x-380\). F: \(y=200x-1000\).
\item Cross at \pounds62, \(11\,400\) yen.
\item \pounds500: E gives \(94\,620\) yen, F gives \(99\,000\) yen; F gives \(4\,380\) yen more. Kenji should use F.
\end{enumerate}
\end{enumerate}

\section*{Part 5: Challenges}

\begin{enumerate}[resume]
\item
\begin{enumerate}
\item G: \(y=1.1x-5.5\), y-intercept \(-5.5\). H: \(y=1.2x-6\), y-intercept \(-6\).
\item Both cross the x-axis at \((5,0)\); \pounds5 gives \texteuro0 at both.
\item Tom is right for any amount over \pounds5. Lines meet only at \((5,0)\); since \(1.2>1.1\), H gives more euros for \(x>5\).
\end{enumerate}

\item Airport at \pounds300 gives \texteuro330. Example new bureau: rate \texteuro1.20 per \pounds1 left, flat fee \pounds25, so \(y=1.2(x-25)=1.2x-30\). At \pounds300 it gives \texteuro330. For amounts more than \pounds300, new bureau is better; for amounts less than \pounds300, airport is better. Other examples: \(r=1.25, f=36\); \(r=1.32, f=50\).

\item
\begin{enumerate}
\item \pounds100 gives \texteuro122.50.
\item Aisha is wrong. Commission is a percentage, so the line passes through the origin. \(y=1.25\times0.98x=1.225x\); gradient \(1.225\).
\item High street: \(y=1.2x-6\). Bureau K: \(y=1.225x\). K is above high street for all positive \(x\); solving gives \(x=-240\), not realistic. K is better for every amount.
\end{enumerate}
\end{enumerate}

\end{document}